\section{Evaluation Practices and Pitfalls}
\label{sec:evaluation}

Cross-domain anomaly detection is unusually easy to evaluate dishonestly by accident. We collect the specific, documented pitfalls a cross-domain study must handle.

\subsection{Metrics: The Point-Adjustment Debate and Its Aftermath}
Point-wise metrics dominated early work and were inflated by point-adjustment (counting any alarm inside an anomalous segment as correct). A line of critique work showed that adjustment can make random detectors look excellent \cite{2109_05257,2308_13068}, proposed unbiased protocols and range-aware metrics \cite{2405_12096,2305_09691}, revisited the ranking behavior of widely used datasets \cite{2603_18985}, and recently stress-tested the post-adjustment evaluation landscape itself \cite{2607_11969}. The field has largely migrated to range-based $F_1$ and VUS, but cross-domain studies add a distinct requirement: since label spaces differ across domains, the \emph{transferable} metrics are score-calibration metrics (coverage, false-alarm rate, interval inflation) plus rank metrics evaluated per target domain---not cross-domain accuracy of a shared classifier.

\subsection{Label and Protocol Provenance}
Benchmark labels arrive as seeded faults, status codes, run-to-failure conventions, or synthetic injections. Each carries a distinct bias: seeded-fault datasets undersample natural degradation; run-to-failure windows are temporally correlated and degradation-ordered; synthetic injections assume separability by construction. Two pathologies we documented empirically \cite{our_scc} deserve special mention. First, \emph{below-chance rankings}: under some labeling protocols, model scores---and even simple feature baselines and zero-shot foundation models---rank anomalous windows below normal ones, because the labeled period differs from normal windows in amplitude or spectral statistics rather than in reconstructability. Second, \emph{alarm-budget illusions}: coverage can look nominal while detection utility is zero. Both pathologies argue for mandatory rankability audits and operational metrics (false alarms per hour, detection delay) alongside calibration numbers.

\subsection{Protocol Desiderata for Cross-Domain Claims}
Synthesizing the critique literature \cite{2109_05257,2409_13053,2511_18739,2405_12096} and our benchmarking experience \cite{our_scc}, we propose the following minimum protocol for cross-domain industrial anomaly detection papers: (i) leave-one-\emph{domain} (ideally leave-one-industry) out, not random splits; (ii) preserve native sampling rates, or evaluate the resampling choice explicitly; (iii) document how every anomaly label was produced and state the protocol's limitations; (iv) report rank metrics with a cheap feature-baseline audit; (v) report a calibration layer with nominal level, achieved coverage, and false-alarm rates separately; (vi) report transfer cost (e.g., interval inflation) so that reliability is priced; and (vii) fix seeds and report seed spread. Most of these are cheap; their absence---not compute---is the main reproducibility gap in the current literature.

\subsection{Classical Baselines and the Reproducibility Baseline}
Comparisons against deep models should include strong classical baselines; several studies find that well-tuned classical detectors remain competitive on industrial streams \cite{2204_01637,2212_11080}, and reconstructed simple baselines (e.g., PCA or feature-Mahalanobis scores) calibrate and rank surprisingly well across domain families. We treat such baselines as mandatory rather than optional in cross-domain evaluations, because they bound how much of the reported gain is attributable to transfer machinery versus score-producer capacity.

\section{Open Problems}
\label{sec:open}

\textbf{Conditional coverage under domain shift.} Marginal coverage is now attainable cross-domain with online conformal methods, but conditional (group- or regime-wise) coverage is not: conditional variants tested to date---Mondrian binning, divergence conditioning---either shrink feedback samples or fail to capture the conditioning that matters \cite{2502_09443,2609_33868}. A principled conditioning variable, and theory connecting measurable diagnostics (exclusion rates, score-distribution drift) to conditional coverage bounds, remain open.

\textbf{Alarm-budget-aware calibration.} The normal-anchored ACI variant resolves the budget allocation empirically, but its exclusion assumption is heuristic; formalizing budget-aware conformal rules---with guarantees that degrade gracefully when anomalies are not separable---is a well-posed theory problem with immediate deployment value in any stream whose anomaly prevalence is unknown \cite{aci2021,2506_15831}.

\textbf{Real process- and mining-industry benchmarks.} The absence of public mining SCADA data and the thinness of semiconductor process benchmarks shape what the field can claim. Fleet-scale corpora \cite{2605_09081} and cross-domain suites \cite{2503_23060} are first steps; what is missing is labeled, native-rate, multi-industry data with documented provenance, plus an evaluation standard \cite{2511_18739} that makes claims comparable.

\textbf{Foundation models versus transfer mechanisms.} Zero-shot foundation models transfer as score producers but do not automatically dominate small domain models on shifted industrial series \cite{freqmatters2025,foundational2025,2505_11349}; understanding \emph{when} scale substitutes for adaptation---and how to combine both---is open, and requires orthogonal evaluations of model capacity and transfer mechanism.

\textbf{Safety of test-time adaptation in monitoring.} TTA objectives can absorb anomalies into the normal model; robust TTA for anomaly detection must actively preserve novelty, not merely reduce reconstruction error \cite{2312_11976,2604_01845}. Criteria and audits for ``adaptation safety'' in monitoring deployments are not yet standardized.

\textbf{Concept drift meets domain shift.} Temporal drift and cross-domain shift are usually studied separately \cite{2506_15831,2412_08435}; real deployments face both, and calibration feedback (alarm-rate control) interacts with drift detectors in ways that are not understood. Joint protocols that vary both axes are a promising empirical direction.

\section{Conclusion}
\label{sec:conclusion}

Cross-domain industrial anomaly detection has matured from per-machine modeling to a literature with a recognizable stack: score producers, transfer mechanisms, calibration layers, and evaluation protocols. This survey argued that the field's center of gravity should shift toward the last two layers. Transfer mechanisms---feature alignment, frequency-aware invariance, test-time adaptation, foundation models, sensor-relation structure---deliver gains that are real but domain-dependent; calibration is where cross-domain deployment is won or lost, because static thresholds fail everywhere, online adaptive conformal methods restore coverage at negligible cost, and the transfer cost is measurable even when it is not predictable. Evaluation, finally, is where the field most often misleads itself: point adjustment is no longer the only trap---below-chance rankings, alarm-budget illusions, and label-protocol bias are systematic. We hope the protocol desiderata and the negative results collected here help the next wave of cross-industry monitoring systems be evaluated the way they will be deployed.
